\chapter{Introduction}
\label{ch:intro}

\section{Background and Motivation}
Assessment is a foundational part of education: it measures what students have learned, guides
teaching, and gives learners feedback on how to improve. A large share of classroom assessment uses
\emph{short-answer questions}, which ask a student to write one to a few sentences that are
then judged for \emph{content correctness} against a reference answer. Grading such answers by hand
is slow and can be inconsistent from one marking session to the next, which limits how often
teachers can give detailed, formative feedback.

\emph{Automatic Short-Answer Grading} (ASAG) studies how to score these answers with a computer. The
field has progressed from early systems based on semantic-similarity and dependency
features~\citep{mohler2011}, through shared tasks that standardised datasets and
evaluation~\citep{dzikovska2013}, to deep neural networks and, most recently, Transformer-based
language models~\citep{vaswani2017attention,devlin2019bert,sung2022survey}. These advances have
markedly improved ASAG for high-resource languages such as English, Arabic, and Chinese.

The benefits, however, have largely \emph{not} reached low-resource languages. \textbf{Khmer}
(\km{ភាសាខ្មែរ}), the official language of Cambodia and spoken by more than 17 million people, is
still treated as low-resource in natural language processing (NLP). Multilingual encoders such as
mBERT, XLM-R, and GTE-multilingual see comparatively little Khmer during pre-training, and the Khmer
script has \emph{no whitespace word boundaries}, so even tokenising a sentence into words is
non-trivial~\citep{buoy2022khmer,khmernltk2020}. Existing Khmer NLP work targets tokenisation,
part-of-speech tagging, and text retrieval, but not educational assessment. At the same time,
artificial intelligence in education is a growing priority for Cambodia~\citep{huot2025cambodia},
which makes a reliable, transparent Khmer grading tool both timely and useful.

A second motivation is \emph{explainability}. A grade is a high-stakes decision: students, teachers,
and parents need to understand \emph{why} an answer received its score, especially if it is to be
trusted or appealed. A teacher will only adopt an automatic grader they can \emph{audit}: one that
visibly rewards the content the rubric cares about rather than spurious surface cues. Crucially, an
explanation is only useful if it is \emph{faithful}: it must reflect what the model actually used,
not merely produce a plausible-looking highlight. Plausible-but-unfaithful explanations are arguably
worse than none, because they create false confidence~\citep{jacovi2020faithfulness}. This thesis
therefore treats explainability as a first-class objective rather than an afterthought.

\section{Problem Statement}
We want to grade Khmer short answers automatically and to \emph{honestly characterise} what is
achievable on a real classroom corpus. Three properties of the available data make this difficult.

\begin{itemize}[leftmargin=1.4em]
  \item \textbf{No benchmark and no baselines.} To our knowledge there is no published Khmer ASAG
  dataset or system to compare against, so the work must build both the benchmark and the reference
  results.
  \item \textbf{Heterogeneous scoring scales.} Each question has its own maximum score, drawn from
  $\{5,6,7,8,10,12,15,20\}$. A raw score of ``8'' therefore means different things on different
  questions, which complicates both modelling and evaluation.
  \item \textbf{A small, single-grader corpus.} The corpus has 1{,}184 answers graded by a
  \emph{single} teacher, so there is limited data and no second grader against whom to measure a
  human agreement ceiling.
\end{itemize}

A further problem, often blurred in ASAG papers, is that two different questions are usually
reported together: \emph{research quality} (how well the model's ordinal scores agree with the
teacher, measured by Quadratic Weighted Kappa) and \emph{deployment quality} (whether the model
predicts the \emph{exact} teacher score in points). This thesis keeps the two separate. Finally, when
a benchmark has few distinct questions, models can memorise question-specific patterns; we must
therefore test how performance holds up on \emph{unseen} questions.

\section{Research Questions}
\label{sec:rqs}
This thesis addresses six research questions.

\begin{description}[leftmargin=2.4em,style=nextline]
  \item[RQ1 -- Model families.] How far can each family of grading model (classical machine
  learning, a recurrent neural network, a Transformer encoder, and a fine-tuned large language
  model) go on a 1{,}184-sample Khmer corpus?
  \item[RQ2 -- Data versus algorithm.] On small-data ASAG, does improving \emph{data quality} move
  results more than swapping the learning algorithm?
  \item[RQ3 -- Research versus deployment.] Does the model that wins the standard research metric
  (QWK) also win the metrics that matter in a classroom (exact-score match and agreement within one
  point)?
  \item[RQ4 -- Cheap levers.] Do two low-cost additions (encoding each question's \emph{maximum
  score} as a feature, and \emph{post-hoc threshold calibration}) improve grading?
  \item[RQ5 -- Explainability (central).] Across the four families, which produces the most
  \emph{faithful} explanations, how well do those explanations align with the teacher's reference
  answer (\emph{plausibility}), and is there an accuracy-versus-explainability trade-off?
  \item[RQ6 -- Robustness.] How large is the \emph{question-leakage} effect, how much does
  performance depend on having seen the question during training, and how well do models grade
  genuinely unseen questions?
\end{description}

\section{Aim and Objectives}
The aim of this thesis is to build the first reproducible, \emph{explainable} ASAG benchmark for
Khmer and to characterise, honestly and with statistical rigour, what current model families can and
cannot do on it. The specific objectives are:

\begin{enumerate}[leftmargin=1.8em]
  \item to construct and document a Khmer short-answer corpus with three curated cleaning variants,
  and a Khmer-aware preprocessing pipeline that addresses script-specific issues;
  \item to implement a single, consistent training-and-evaluation pipeline that compares four model
  families as bounded ordinal regression;
  \item to design a \emph{model-agnostic explainability protocol} that measures explanation
  faithfulness (using ERASER comprehensiveness/sufficiency and AOPC) and plausibility, and to apply
  it across all families;
  \item to evaluate every model with field-standard metrics, reported consistently across three
  dataset variants, and to quantify the question-leakage effect;
  and
  \item to realise the explainability contribution as a live, teacher-facing web prototype.
\end{enumerate}

\section{Significance of the Study}
This study makes the first contribution of automated short-answer grading technology to the Khmer
language, providing a reproducible benchmark, baselines across four model families, and a released
dataset and codebase that future Khmer NLP work can build on. Methodologically, it contributes a
cross-family \emph{faithfulness} protocol that measures, rather than assumes, that explanations
reflect the model's reasoning, together with an honest \emph{question-leakage} analysis that
quantifies how much small-question-pool ASAG scores are inflated. Practically, the live prototype demonstrates how
a transparent grading assistant could give Cambodian teachers faster, auditable feedback. The thesis
is also written to be \emph{honest by construction}: every headline number is tied to a verified
computed result, results are reported on both train and test across three dataset variants, and the
limits of the work are stated plainly.

\section{Scope and Limitations}
The thesis focuses on short-answer questions (typically one to a few sentences) in Khmer, drawn from
secondary-school subjects (Biology, History, Geography, and Earth Science). It studies
\emph{content} scoring against a reference answer, not the grading of essays for style or grammar. A
single teacher graded every answer, so the work reports agreement \emph{with that grader} rather
than against a measured human ceiling; obtaining a second annotator was not possible within the
scope of the study. The most computationally heavy models (the Transformer encoders and the
fine-tuned LLM) require a GPU and were run on dedicated hardware; the classical and recurrent
baselines run on a CPU. These limitations are revisited in Chapter~\ref{ch:limitations}.

\section{Thesis Organisation}
The remainder of the thesis is organised as follows. Chapter~\ref{ch:literature} reviews the
literature on ASAG, Transformer models, explainable AI and faithfulness, cross-lingual and
low-resource scoring, and Khmer NLP, and identifies the research gap. Chapter~\ref{ch:methodology}
describes the research design, the dataset and preprocessing, the four model pillars, the
explainability protocol, and the evaluation methodology. Chapter~\ref{ch:experiments} reports the
experimental setup, the experiment grid, hyperparameter tuning, and the live prototype.
Chapter~\ref{ch:results} presents and discusses the results, per-pillar performance with confidence
intervals, the research-versus-deployment trade-off, the explainability study, and the leakage
analysis. Chapter~\ref{ch:limitations} states limitations and future work, and
Chapter~\ref{ch:conclusion} concludes.
